\documentclass[12pt,a4paper]{article}

\usepackage[T2A]{fontenc}
\usepackage[utf8]{inputenc}
\usepackage[russian]{babel}
\usepackage{geometry}
\geometry{left=30mm,right=20mm,top=20mm,bottom=20mm}
\usepackage{setspace}
\onehalfspacing
\usepackage{graphicx}
\usepackage{amsmath,amssymb}
\usepackage{float}
\usepackage{caption}
\usepackage{enumitem}
\usepackage{hyperref}
\usepackage{booktabs}
\usepackage{array}
\usepackage{indentfirst}
\usepackage{titlesec}
\usepackage{tocloft}

% Заголовки РАЗДЕЛОВ по центру
\titleformat{\section}
  {\centering\large\bfseries}
  {\thesection}
  {1em}
  {}
\titleformat{\subsection}
  {\normalsize\bfseries}
  {\thesubsection}
  {1em}
  {}
\titleformat{\subsubsection}
  {\normalsize\bfseries}
  {\thesubsubsection}
  {1em}
  {}

% Списки: длинное тире
\setlist[itemize]{label=---, leftmargin=*}
\setlist[enumerate]{leftmargin=*}

% Содержание — заголовок СЛЕВА
\renewcommand{\contentsname}{Содержание}

\numberwithin{equation}{section}

% Абзацный отступ
\setlength{\parindent}{1.25cm}
